% GENERATED FILE. Edit canonical lessons, then run npm run generate:books.
% canonical-source-hash: fnv1a64:01f07be026276f6b
% canonical-lessons: AR-C319-aza, AR-C319-azza, AR-C319-ratha, AR-C319-naa, AR-C319-hidad

\chapter{Mourning}
\label{ch:ar-mourning}
\hlchaptermodality{319}

\begin{chapteropening}
\textbf{By the end of this chapter:} \emph{I can say condolence, to console, to lament, to announce a death and mourning.}

Complete the last lesson of chapter 319: I can say condolence, to console, to lament, to announce a death and mourning.
\end{chapteropening}

\section[\texorpdfstring{\ar{عزاء} (ʿazāʾ)}{عزاء (ʿazāʾ)}]{\ar{عزاء} (ʿazāʾ) — condolence; a gathering of mourners}

\label{lesson:AR-C319-aza}



\begin{quote}
Before the new one: say the Arabic for a temple, then the Arabic for a pyramid.
\end{quote}

\subsection*{You'll want to know: \ar{عزاء}}
\textbf{\ar{عزاء}} — \emph{ʿazāʾ} — \textquotedblleft{}condolence; a gathering of mourners\textquotedblright{}.

\begin{quote}
\ar{ذهبنا} \ar{إلى} \ar{بيت} \ar{العزاء} \ar{مساء} \ar{أمس}.
\end{quote}

\begin{quote}
\emph{dhahabnā ilā bayt al-ʿazāʾ masāʾa ams.}
\end{quote}

\begin{quote}
We went to the house of mourning yesterday evening.
\end{quote}

\ar{بيت} \ar{العزاء} is the home where the family receives visitors after a death.

\begin{cousinweb}
\textbf{\ar{عزاء}} (\emph{ʿazāʾ}) is built on the native Arabic root \emph{ʿ-z-y} (Lane), where \emph{ʿazāʾ} is patience in the face of loss. To offer it to a family is to wish them that patience.
\end{cousinweb}

\begin{grammarlens}[title={What you've built}]
One more word for telling what happened.
\end{grammarlens}

\subsection*{Your turn}
\begin{itemize}
\raggedright
  \item \emph{Say it:} \emph{ʿazāʾ}
  \item \emph{Say it:} \emph{ʿazāʾ}, once more
  \item \emph{Say it:} \emph{dhahabnā ilā bayt al-ʿazāʾ masāʾa ams.} — We went to the house of mourning yesterday evening.
  \item \emph{Recall:} say the Arabic for a temple, then the Arabic for a pyramid, then say \emph{ʿazāʾ} again
\end{itemize}

\begin{tcolorbox}[breakable,colback=teal!4,colframe=teal!35!black,title={Before you move on}]
What does \ar{عزاء} mean? (\textquotedblleft{}Condolence; a gathering of mourners\textquotedblright{}.) Say it once more.
\end{tcolorbox}
\section[\texorpdfstring{\ar{عزى} (ʿazzā)}{عزى (ʿazzā)}]{\ar{عزى} (ʿazzā) — to console, to offer condolences}

\label{lesson:AR-C319-azza}



\begin{quote}
Before the new one: say the Arabic for a pyramid, then the Arabic for condolence.
\end{quote}

\subsection*{You'll want to know: \ar{عزى}}
\textbf{\ar{عزى}} — \emph{ʿazzā} — \textquotedblleft{}to console, to offer condolences\textquotedblright{}.

\begin{quote}
\ar{عزى} \ar{الجيران} \ar{الأسرة} \ar{بعد} \ar{الجنازة}.
\end{quote}

\begin{quote}
\emph{ʿazzā al-jīrān al-usra baʿda al-janāza.}
\end{quote}

\begin{quote}
The neighbours offered the family their condolences after the funeral.
\end{quote}

Present: \ar{يعزي}, \emph{yuʿazzī}.

\ar{عزى} someone \ar{في} the person who died: \ar{عزاه} \ar{في} \ar{أبيه}.

\begin{cousinweb}
\textbf{\ar{عزى}} (\emph{ʿazzā}) is built on the native Arabic root \emph{ʿ-z-y} (Lane), patience under loss. It is Form II, \textquotedblleft{}to bid someone be patient\textquotedblright{}: to console.
\end{cousinweb}

\begin{grammarlens}[title={What you've built}]
One more word for telling what happened.
\end{grammarlens}

\subsection*{Your turn}
\begin{itemize}
\raggedright
  \item \emph{Say it:} \emph{ʿazzā}
  \item \emph{Say it:} \emph{ʿazzā}, once more
  \item \emph{Say it:} \emph{ʿazzā al-jīrān al-usra baʿda al-janāza.} — The neighbours offered the family their condolences after the funeral.
  \item \emph{Recall:} say the Arabic for a pyramid, then the Arabic for condolence, then say \emph{ʿazzā} again
\end{itemize}

\begin{tcolorbox}[breakable,colback=teal!4,colframe=teal!35!black,title={Before you move on}]
What does \ar{عزى} mean? (\textquotedblleft{}To console, to offer condolences\textquotedblright{}.) Say it once more.
\end{tcolorbox}
\section[\texorpdfstring{\ar{رثى} (rathā)}{رثى (rathā)}]{\ar{رثى} (rathā) — to lament, to mourn in verse}

\label{lesson:AR-C319-ratha}



\begin{quote}
Before the new one: say the Arabic for condolence, then the Arabic for to console.
\end{quote}

\subsection*{You'll want to know: \ar{رثى}}
\textbf{\ar{رثى}} — \emph{rathā} — \textquotedblleft{}to lament, to mourn in verse\textquotedblright{}.

\begin{quote}
\ar{رثى} \ar{الشاعر} \ar{صديقه} \ar{في} \ar{قصيدة} \ar{طويلة}.
\end{quote}

\begin{quote}
\emph{rathā ash-shāʿir ṣadīqahu fī qaṣīda ṭawīla.}
\end{quote}

\begin{quote}
The poet mourned his friend in a long poem.
\end{quote}

Present: \ar{يرثي}, \emph{yarthī}.

A poem of mourning is \ar{رثاء}, an elegy.

\begin{cousinweb}
\textbf{\ar{رثى}} (\emph{rathā}) is built on the native Arabic root \emph{r-th-y} (Lane), \textquotedblleft{}to weep for the dead and recount his good deeds\textquotedblright{}; \emph{rithāʾ}, the elegy, is one of the oldest kinds of Arabic poem.
\end{cousinweb}

\begin{grammarlens}[title={What you've built}]
One more word for telling what happened.
\end{grammarlens}

\subsection*{Your turn}
\begin{itemize}
\raggedright
  \item \emph{Say it:} \emph{rathā}
  \item \emph{Say it:} \emph{rathā}, once more
  \item \emph{Say it:} \emph{rathā ash-shāʿir ṣadīqahu fī qaṣīda ṭawīla.} — The poet mourned his friend in a long poem.
  \item \emph{Recall:} say the Arabic for condolence, then the Arabic for to console, then say \emph{rathā} again
\end{itemize}

\begin{tcolorbox}[breakable,colback=teal!4,colframe=teal!35!black,title={Before you move on}]
What does \ar{رثى} mean? (\textquotedblleft{}To lament, to mourn in verse\textquotedblright{}.) Say it once more.
\end{tcolorbox}
\section[\texorpdfstring{\ar{نعى} (naʿā)}{نعى (naʿā)}]{\ar{نعى} (naʿā) — to announce a death}

\label{lesson:AR-C319-naa}



\begin{quote}
Before the new one: say the Arabic for to console, then the Arabic for to lament.
\end{quote}

\subsection*{You'll want to know: \ar{نعى}}
\textbf{\ar{نعى}} — \emph{naʿā} — \textquotedblleft{}to announce a death\textquotedblright{}.

\begin{quote}
\ar{نعت} \ar{الصحيفة} \ar{الكاتب} \ar{الكبير}.
\end{quote}

\begin{quote}
\emph{naʿat aṣ-ṣaḥīfa al-kātib al-kabīr.}
\end{quote}

\begin{quote}
The newspaper announced the death of the great writer.
\end{quote}

Present: \ar{ينعى}, \emph{yanʿā}.

A death notice in the paper is \ar{نعي}.

\begin{cousinweb}
\textbf{\ar{نعى}} (\emph{naʿā}) is built on the native Arabic root \emph{n-ʿ-y} (Lane), \textquotedblleft{}to bring the news of someone's death\textquotedblright{}; \emph{naʿy} is that news.
\end{cousinweb}

\begin{grammarlens}[title={What you've built}]
One more word for telling what happened.
\end{grammarlens}

\subsection*{Your turn}
\begin{itemize}
\raggedright
  \item \emph{Say it:} \emph{naʿā}
  \item \emph{Say it:} \emph{naʿā}, once more
  \item \emph{Say it:} \emph{naʿat aṣ-ṣaḥīfa al-kātib al-kabīr.} — The newspaper announced the death of the great writer.
  \item \emph{Recall:} say the Arabic for to console, then the Arabic for to lament, then say \emph{naʿā} again
\end{itemize}

\begin{tcolorbox}[breakable,colback=teal!4,colframe=teal!35!black,title={Before you move on}]
What does \ar{نعى} mean? (\textquotedblleft{}To announce a death\textquotedblright{}.) Say it once more.
\end{tcolorbox}
\section[\texorpdfstring{\ar{حداد} (ḥidād)}{حداد (ḥidād)}]{\ar{حداد} (ḥidād) — mourning, a period of mourning}

\label{lesson:AR-C319-hidad}



\begin{quote}
Before the new one: say the Arabic for to lament, then the Arabic for to announce a death.
\end{quote}

\subsection*{You'll want to know: \ar{حداد}}
\textbf{\ar{حداد}} — \emph{ḥidād} — \textquotedblleft{}mourning, a period of mourning\textquotedblright{}.

\begin{quote}
\ar{لبست} \ar{الأسرة} \ar{الأسود} \ar{في} \ar{أيام} \ar{الحداد}.
\end{quote}

\begin{quote}
\emph{labisat al-usra al-aswad fī ayyām al-ḥidād.}
\end{quote}

\begin{quote}
The family wore black in the days of mourning.
\end{quote}

\ar{أعلنت} \ar{الحكومة} \ar{الحداد} \ar{ثلاثة} \ar{أيام}: three days of national mourning.

\begin{cousinweb}
\textbf{\ar{حداد}} (\emph{ḥidād}) is built on the native Arabic root \emph{ḥ-d-d} (Lane), the root of \emph{ḥadd}, a limit, and of \emph{ḥadīd}, iron. Lane gives \emph{ḥidād} as a widow's holding back from finery and perfume: mourning.
\end{cousinweb}

\begin{grammarlens}[title={What you've built}]
One more word for telling what happened.
\end{grammarlens}

\subsection*{Your turn}
\begin{itemize}
\raggedright
  \item \emph{Say it:} \emph{ḥidād}
  \item \emph{Say it:} \emph{ḥidād}, once more
  \item \emph{Say it:} \emph{labisat al-usra al-aswad fī ayyām al-ḥidād.} — The family wore black in the days of mourning.
  \item \emph{Recall:} say the Arabic for to lament, then the Arabic for to announce a death, then say \emph{ḥidād} again
\end{itemize}

\begin{tcolorbox}[breakable,colback=teal!4,colframe=teal!35!black,title={Before you move on}]
What does \ar{حداد} mean? (\textquotedblleft{}Mourning, a period of mourning\textquotedblright{}.) Say it once more.
\end{tcolorbox}
